\section{结论}
\label{sec:conclusion}

本文围绕“把一枚教学用运载火箭送入 600\,km 近圆轨道”这一目标，完成了建模、数值求解、
轨迹优化、闭环控制与约束敏感性的完整研究，主要结论如下。

\begin{enumerate}
\item \textbf{模型与数值层面。}在平面极坐标变质量动力学框架内，用同一组方程服务
      规划、预测与实际对象；采用三档积分精度，并明确只有高精度档的结果允许写入结论。
      自洽性核对显示推力律 $T=cq$ 精确成立、冲量与质量关系相对误差小于 $10^{-8}$；
      步长细化表明终端半径的数值不确定度约 3.18\,m，因此\textbf{可信精度为“数米级”}。

\item \textbf{规划层面。}10 维两段燃烧参数化 + PSO 全局搜索 + SQP 局部精修，
      在 117.5\,s 内得到内点解：关机时刻 577.305\,s，终端高度误差 8.281\,m、
      径向速度 0.0446\,m/s、切向速度误差 0.020\,m/s，瞬时密切轨道偏心率
      $8.69\times10^{-6}$；消耗推进剂 52.69\,t（预算 58.8\,t），终端质量 7.31\,t。
      方案把圆轨化推迟到 599.60\,km 的远地点，结构接近霍曼转移。

\item \textbf{约束敏感性。}紧约束解（偏角 $\pm10^\circ$、滑行 $\le250$\,s）的三个变量
      贴在上界；仅放宽这三处，同一终端条件下推进剂少用 1.08\,t（$-2.0\%$）、
      峰值过载由 15.9\,g 降到 13.1\,g，代价是动压峰值由 47.8\,kPa 升到 68.6\,kPa、
      飞行时间增加 112\,s。这说明\textbf{原解确实被盒约束截断}，也说明
      \textbf{多出的自由度被用来换燃料而不是换精度}（终端残差由 0.265\,m 变为 8.281\,m）。

\item \textbf{加速层面。}GPU 批量定步长 RK4 把 96 条轨迹的批次时间从 9.834\,s
      降到 0.324\,s（30.4 倍），但目标值相对差 2.03\%，说明它只适合做候选排序，
      不能替代高精度复核。

\item \textbf{验证层面。}关机后仅受引力与阻力推进一个标称周期 5792.33\,s，
      高度维持在 599.992--600.113\,km，极差 121\,m 与密切偏心率 $8.69\times10^{-6}$
      自洽。这支持“瞬时近圆注入”的判断，但不支持“长期保持轨道”的更强结论。

\item \textbf{控制层面。}在独立扰动实现下，LTV-MPC 把位置跟踪 RMS 从开环的
      31.5\,km 降到 4.21\,m，速度误差从 107.96\,m/s 降到 0.089\,m/s；
      286 次 QP 全部收敛、平均 2.24\,ms，最大修正仅 $0.276^\circ$ 与 0.71\,kg/s，
      说明闭环工作在“小偏差”区间，与线性化假设自洽。

\item \textbf{局限层面。}结果对应单次扰动试验、理想全状态反馈与简化大气模型；
      动压 68.6\,kPa 与过载 13.1\,g 都超出工程常用范围，暴露了“固定流率、不节流”
      假设的不合理性。因此本文不声称全局最优、统计鲁棒性或严格闭环稳定性。
\end{enumerate}

从方法上看，本文最值得保留的经验不是某个算法，而是\textbf{把“模型—精度—算法”
三者的关系讲清楚}：模型要统一，精度要分级，算法要分工——而分工本身取决于
问题处在可行域的什么位置。紧约束时 SQP 是必需品；放宽边界后它几乎无事可做。
只有同时报告“解”和“边界”，结论才有意义。

后续工作有三个自然的延伸方向：一是把 $q_d\le50$\,kPa 与 $T/m\le6g$ 写进约束，
以燃料最小为目标重跑，回答“放宽边界省下的燃料在工程约束下还剩多少”；
二是做多种子蒙特卡洛与边界参数扫描，把“单次成功”升级为“成功率”与
“燃料—动压前沿”；三是把 LTV-MPC 扩展为管状 MPC，在保留凸性的同时获得鲁棒性保证。
